\section{Sort + search tổng hợp [Trung bình]}
\subsection{Phân tích \& Thiết kế (M0)}

\begin{itemize}

    \item \textbf{Input:}

    File \texttt{records.csv} chứa:

\begin{lstlisting}
id,score
\end{lstlisting}

    Ví dụ:

\begin{lstlisting}
A,9.5
B,7.0
C,9.5
\end{lstlisting}

    File \texttt{queries.txt} chứa:

    \begin{itemize}
        \item Mỗi dòng là một \texttt{id} cần tra cứu
    \end{itemize}

    \item \textbf{Xử lý:}

    \begin{itemize}

        \item \textbf{Case biên:}

        \begin{itemize}
            \item File records rỗng
            \[
                \Rightarrow \text{mọi query đều NOT\_FOUND}
            \]

            \item Query không tồn tại
            \[
                \Rightarrow \texttt{NOT\_FOUND}
            \]

            \item Các record có cùng score
            \[
                \Rightarrow \text{giữ nguyên thứ tự ban đầu}
            \]
        \end{itemize}

        \item \textbf{Module 1 (Đọc file records):}

        \begin{itemize}
            \item Hàm \texttt{readRecords()}
        \end{itemize}

        \item \textbf{Module 2 (Đọc query):}

        \begin{itemize}
            \item Hàm \texttt{readQueries()}
        \end{itemize}

        \item \textbf{Module 3 (Sort records):}

        \begin{itemize}
            \item Hàm \texttt{sortRecords()}
            \item Dùng:
\begin{lstlisting}
stable_sort()
\end{lstlisting}

            \item Sort theo score giảm dần
        \end{itemize}

        \item \textbf{Module 4 (Search):}

        \begin{itemize}
            \item Hàm \texttt{findRecord()}
            \item Tìm:
            \begin{itemize}
                \item rank (1-based)
                \item id
                \item score
            \end{itemize}
        \end{itemize}

        \item \textbf{Module 5 (Ghi file):}

        \begin{itemize}
            \item Hàm \texttt{writeOutput()}
        \end{itemize}

    \end{itemize}

    \item \textbf{Output:}

    File \texttt{output.txt}

\begin{lstlisting}
rank,id,score
NOT_FOUND
...
\end{lstlisting}

    Dùng:
\begin{lstlisting}
setprecision(6)
\end{lstlisting}

\end{itemize}

\subsection{Mã nguồn C++11 (Tách module)}

\begin{lstlisting}
#include <iostream>
#include <fstream>
#include <sstream>
#include <vector>
#include <algorithm>
#include <iomanip>

using namespace std;

struct Record {

    string id;

    double score;
};

bool readRecords(
    const string& path,
    vector<Record>& records
) {
    ifstream fin(path);

    if (!fin) return false;

    string line;

    while (getline(fin, line)) {

        if (line.empty()) continue;

        stringstream ss(line);

        string id, scoreStr;

        getline(ss, id, ',');
        getline(ss, scoreStr);

        Record r;

        r.id = id;

        r.score = stod(scoreStr);

        records.push_back(r);
    }

    return true;
}

bool readQueries(
    const string& path,
    vector<string>& queries
) {
    ifstream fin(path);

    if (!fin) return false;

    string id;

    while (getline(fin, id)) {

        if (!id.empty()) {
            queries.push_back(id);
        }
    }

    return true;
}

void sortRecords(
    vector<Record>& records
) {
    stable_sort(
        records.begin(),
        records.end(),

        [](const Record& a,
           const Record& b) {

            return a.score > b.score;
        }
    );
}

int findRecord(
    const vector<Record>& records,
    const string& id
) {
    for (size_t i = 0; i < records.size(); ++i) {

        if (records[i].id == id) {
            return i;
        }
    }

    return -1;
}

bool writeOutput(
    const string& path,
    const vector<Record>& records,
    const vector<string>& queries
) {
    ofstream fout(path);

    if (!fout) return false;

    fout << fixed << setprecision(6);

    for (size_t i = 0; i < queries.size(); ++i) {

        int pos =
            findRecord(records, queries[i]);

        if (pos == -1) {

            fout << "NOT_FOUND\n";
        }
        else {

            fout << pos + 1 << ","
                 << records[pos].id << ","
                 << records[pos].score
                 << "\n";
        }
    }

    return true;
}

int main() {

    vector<Record> records;

    if (!readRecords(
        "records.csv",
        records
    )) {
        cerr << "Loi mo file records!\n";
        return 1;
    }

    vector<string> queries;

    if (!readQueries(
        "queries.txt",
        queries
    )) {
        cerr << "Loi mo file queries!\n";
        return 1;
    }

    sortRecords(records);

    if (!writeOutput(
        "output.txt",
        records,
        queries
    )) {
        cerr << "Loi mo file output!\n";
        return 1;
    }

    return 0;
}
\end{lstlisting}

\subsection{Kế hoạch kiểm thử (Test Cases)}

\textbf{Test Case 1 (Thông thường)}

\textbf{Input (records.csv):}

\begin{lstlisting}
A,9.5
B,7.0
C,9.5
D,6.0
\end{lstlisting}

\textbf{Input (queries.txt):}

\begin{lstlisting}
C
B
D
\end{lstlisting}

\textbf{Expected Output (output.txt):}

\begin{lstlisting}
2,C,9.500000
3,B,7.000000
4,D,6.000000
\end{lstlisting}

\vspace{0.5cm}

\textbf{Test Case 2 (Không tìm thấy)}

\textbf{Input (records.csv):}

\begin{lstlisting}
A,10
B,20
\end{lstlisting}

\textbf{Input (queries.txt):}

\begin{lstlisting}
C
A
\end{lstlisting}

\textbf{Expected Output (output.txt):}

\begin{lstlisting}
NOT_FOUND
2,A,10.000000
\end{lstlisting}

\vspace{0.5cm}

\textbf{Test Case 3 (File records rỗng)}

\textbf{Input (records.csv):}

\begin{lstlisting}

\end{lstlisting}

\textbf{Input (queries.txt):}

\begin{lstlisting}
A
B
\end{lstlisting}

\textbf{Expected Output (output.txt):}

\begin{lstlisting}
NOT_FOUND
NOT_FOUND
\end{lstlisting}